% Chapter 7-8: Abelian Categories and Derived Categories
% Continuation -- no preamble

\chapter{Abelian Categories and Exact Sequences}\label{ch:abelian}
\minitoc

\section{Preadditive and Additive Categories}\label{sec:preadditive}

The passage from general category theory to homological algebra requires
enriching the categorical framework with algebraic structure on morphism sets.
This section introduces the hierarchy of additive structures on categories.

\begin{definition}[Preadditive category]\label{def:preadditive-cat}
\index{category!preadditive}\index{preadditive category}
A \textbf{preadditive category} (or \textbf{$\Ab$-category}) is a category $\mathcal{A}$
such that:
\begin{enumerate}
    \item For every pair of objects $A, B \in \Ob(\mathcal{A})$, the set $\Hom_{\mathcal{A}}(A,B)$
    carries the structure of an abelian group, with addition denoted $+$ and zero element
    denoted $0_{A,B}$.
    \item Composition is bilinear: for all morphisms $f, f' \colon A \to B$ and
    $g, g' \colon B \to C$,
    \[
        g \circ (f + f') = g \circ f + g \circ f', \qquad
        (g + g') \circ f = g \circ f + g' \circ f.
    \]
\end{enumerate}
Equivalently, $\mathcal{A}$ is a category enriched over $(\Ab, \otimes_{\mathbb{Z}}, \mathbb{Z})$.
\end{definition}

\begin{remark}\label{rem:preadditive-zero}
\index{zero morphism}
In a preadditive category, for any objects $A, B$, the zero element $0_{A,B} \in \Hom(A,B)$
satisfies $0_{B,C} \circ f = 0_{A,C}$ and $g \circ 0_{A,B} = 0_{A,C}$ for all
$f \colon A \to B$ and $g \colon B \to C$. Thus the zero morphisms form a compatible system.
If $\mathcal{A}$ has a zero object $0$, then $0_{A,B}$ is the composite $A \to 0 \to B$.
\end{remark}

\begin{example}\label{ex:preadditive-examples}
\index{Ab@$\Ab$!as preadditive category}\index{Mod@$\Mod$!as preadditive category}
\begin{enumerate}
    \item The category $\Ab$ of abelian groups is preadditive, with pointwise addition
    of homomorphisms.
    \item For any ring $R$, the category $\Mod_R$ of right $R$-modules is preadditive.
    \item The category $\Set$ is \emph{not} preadditive: there is no natural abelian group
    structure on $\Hom_{\Set}(X,Y)$ for general sets $X, Y$.
    \item Any ring $R$ can be viewed as a preadditive category with a single object $*$,
    where $\Hom(*, *) = R$.
\end{enumerate}
\end{example}

\begin{definition}[Additive functor]\label{def:additive-functor}
\index{functor!additive}\index{additive functor}
Let $\mathcal{A}, \mathcal{B}$ be preadditive categories. A functor
$F \colon \mathcal{A} \to \mathcal{B}$ is \textbf{additive} if for all objects $A, B$
the map
\[
    F_{A,B} \colon \Hom_{\mathcal{A}}(A,B) \longrightarrow \Hom_{\mathcal{B}}(FA, FB)
\]
is a group homomorphism: $F(f + g) = F(f) + F(g)$.
\end{definition}

\begin{definition}[Biproduct]\label{def:biproduct}
\index{biproduct}\index{direct sum!in additive category}
Let $\mathcal{A}$ be a preadditive category and let $A_1, \ldots, A_n$ be objects.
A \textbf{biproduct} of $A_1, \ldots, A_n$ is an object $A_1 \oplus \cdots \oplus A_n$
together with morphisms
\[
    \iota_k \colon A_k \to A_1 \oplus \cdots \oplus A_n, \qquad
    \pi_k \colon A_1 \oplus \cdots \oplus A_n \to A_k, \qquad 1 \leq k \leq n,
\]
satisfying:
\begin{enumerate}
    \item $\pi_k \circ \iota_j = \delta_{jk} \cdot \id_{A_k}$ for all $j, k$,
    \item $\sum_{k=1}^n \iota_k \circ \pi_k = \id_{A_1 \oplus \cdots \oplus A_n}$.
\end{enumerate}
\end{definition}

\begin{proposition}\label{prop:biproduct-universal}
\index{biproduct!universal property}
A biproduct $(A_1 \oplus \cdots \oplus A_n, \iota_k, \pi_k)$ is simultaneously
a product and a coproduct. More precisely:
\begin{enumerate}
    \item $(A_1 \oplus \cdots \oplus A_n, (\pi_k))$ is a product of $A_1, \ldots, A_n$.
    \item $(A_1 \oplus \cdots \oplus A_n, (\iota_k))$ is a coproduct of $A_1, \ldots, A_n$.
\end{enumerate}
\end{proposition}

\begin{proof}
We prove the product property; the coproduct property is dual. Given morphisms
$f_k \colon B \to A_k$ for $1 \leq k \leq n$, define
$f = \sum_{k=1}^n \iota_k \circ f_k \colon B \to A_1 \oplus \cdots \oplus A_n$.
Then
\[
    \pi_j \circ f = \sum_{k=1}^n \pi_j \circ \iota_k \circ f_k
    = \sum_{k=1}^n \delta_{jk} f_k = f_j.
\]
For uniqueness, if $g \colon B \to A_1 \oplus \cdots \oplus A_n$ satisfies
$\pi_k \circ g = f_k$ for all $k$, then
\[
    g = \Bigl(\sum_{k=1}^n \iota_k \circ \pi_k\Bigr) \circ g
    = \sum_{k=1}^n \iota_k \circ f_k = f. \qedhere
\]
\end{proof}

\begin{definition}[Additive category]\label{def:additive-cat}
\index{category!additive}\index{additive category}
An \textbf{additive category} is a preadditive category $\mathcal{A}$ that:
\begin{enumerate}
    \item has a zero object $0$,
    \item has finite biproducts: for any objects $A, B \in \Ob(\mathcal{A})$,
    the biproduct $A \oplus B$ exists.
\end{enumerate}
\end{definition}

\begin{proposition}\label{prop:additive-criterion}
\index{additive category!characterization}
Let $\mathcal{A}$ be a preadditive category with a zero object and finite products
(or finite coproducts). Then $\mathcal{A}$ is additive.
\end{proposition}

\begin{proof}
Suppose $\mathcal{A}$ has finite products. Let $A_1, A_2 \in \Ob(\mathcal{A})$,
and let $(A_1 \times A_2, \pi_1, \pi_2)$ be their product. Define
$\iota_1 \colon A_1 \to A_1 \times A_2$ as the unique morphism with
$\pi_1 \circ \iota_1 = \id_{A_1}$ and $\pi_2 \circ \iota_1 = 0$,
and similarly for $\iota_2$. One checks that $\iota_1 \circ \pi_1 + \iota_2 \circ \pi_2 = \id$
using the universal property and the fact that
\[
    \pi_j \circ (\iota_1 \circ \pi_1 + \iota_2 \circ \pi_2) = \delta_{1j}\pi_1 + \delta_{2j}\pi_2 = \pi_j
\]
for $j = 1, 2$. Thus $(A_1 \times A_2, \iota_k, \pi_k)$ is a biproduct.
\end{proof}

\begin{example}\label{ex:biproduct-matrix}
\index{biproduct!matrix description}
In $\Mod_R$, the biproduct $M \oplus N$ is the usual direct sum. The morphism set
$\Hom(M_1 \oplus M_2, N_1 \oplus N_2)$ can be identified with $2 \times 2$ matrices:
\[
    \begin{pmatrix} f_{11} & f_{12} \\ f_{21} & f_{22} \end{pmatrix},
    \qquad f_{ij} \colon M_j \to N_i,
\]
with composition given by matrix multiplication. This matrix calculus is a key
feature of additive categories.
\end{example}


\section{Kernels, Cokernels, and Abelian Categories}\label{sec:abelian-def}

\begin{definition}[Kernel and cokernel]\label{def:kernel-cokernel}
\index{kernel!in additive category}\index{cokernel}
Let $f \colon A \to B$ be a morphism in a preadditive category $\mathcal{A}$.
\begin{enumerate}
    \item A \textbf{kernel} of $f$ is an equalizer of $f$ and $0_{A,B}$: it is a morphism
    $k \colon K \to A$ such that $f \circ k = 0$ and for every $g \colon X \to A$ with
    $f \circ g = 0$, there exists a unique $\bar{g} \colon X \to K$ with $k \circ \bar{g} = g$.
    % Diagram
    \[
    \begin{tikzcd}
        X \arrow[rd, "g"'] \arrow[rrd, "0", bend left=20] \arrow[d, dashed, "\exists!\bar{g}"'] \\
        K \arrow[r, "k"'] & A \arrow[r, "f"'] & B
    \end{tikzcd}
    \]
    We write $\Ker(f) = (K, k)$ or simply $\Ker(f) = K$.
    \item A \textbf{cokernel} of $f$ is a coequalizer of $f$ and $0_{A,B}$: it is a morphism
    $c \colon B \to C$ such that $c \circ f = 0$ and for every $h \colon B \to Y$ with
    $h \circ f = 0$, there exists a unique $\bar{h} \colon C \to Y$ with $\bar{h} \circ c = h$.
    We write $\coker(f) = (C, c)$.
\end{enumerate}
\end{definition}

\begin{definition}[Image and coimage]\label{def:image-coimage}
\index{image!categorical}\index{coimage}
Let $f \colon A \to B$ be a morphism in a preadditive category.
\begin{enumerate}
    \item The \textbf{image} of $f$ is $\im(f) = \Ker(\coker(f))$.
    \item The \textbf{coimage} of $f$ is $\operatorname{coim}(f) = \coker(\Ker(f))$.
\end{enumerate}
There is a canonical morphism $\overline{f} \colon \operatorname{coim}(f) \to \im(f)$
induced by the universal properties.
\end{definition}

\begin{definition}[Abelian category]\label{def:abelian-cat}
\index{abelian category}\index{category!abelian}
An \textbf{abelian category} is an additive category $\mathcal{A}$ satisfying:
\begin{enumerate}
    \item[(AB1)] Every morphism has a kernel and a cokernel.
    \item[(AB2)] For every morphism $f$, the canonical morphism
    $\overline{f} \colon \operatorname{coim}(f) \to \im(f)$ is an isomorphism.
\end{enumerate}
\end{definition}

\begin{remark}\label{rem:abelian-equiv}
\index{abelian category!equivalent formulations}
Condition (AB2) can equivalently be stated: every monomorphism is a kernel and
every epimorphism is a cokernel. That is, every mono $m \colon K \hookrightarrow A$ is
the kernel of some morphism out of $A$, and dually for epimorphisms.
\end{remark}

\begin{example}\label{ex:abelian-examples}
\begin{enumerate}
    \item\index{Ab@$\Ab$!as abelian category}
    $\Ab$ is an abelian category. Kernels and cokernels are the usual ones.
    \item\index{Mod@$\Mod$!as abelian category}
    $\Mod_R$ is abelian for any ring $R$. This is the prototypical example.
    \item\index{sheaf!category of sheaves}
    For a topological space $X$, the category $\operatorname{Sh}(X, \Ab)$ of sheaves of
    abelian groups on $X$ is abelian.
    \item The category of finitely generated abelian groups is abelian.
    \item\index{category!not abelian}
    The category of free abelian groups is additive but \emph{not} abelian: the cokernel
    of $\mathbb{Z} \xrightarrow{2} \mathbb{Z}$ would be $\mathbb{Z}/2\mathbb{Z}$, which
    is not free.
\end{enumerate}
\end{example}

\begin{definition}[Grothendieck axioms]\label{def:grothendieck-axioms}
\index{Grothendieck axioms}\index{AB3}\index{AB4}\index{AB5}\index{AB3*}\index{AB4*}\index{AB5*}
Let $\mathcal{A}$ be an abelian category. Grothendieck introduced additional axioms:
\begin{enumerate}
    \item[(AB3)] $\mathcal{A}$ has arbitrary (set-indexed) coproducts.
    \item[(AB3*)] $\mathcal{A}$ has arbitrary (set-indexed) products.
    \item[(AB4)] (AB3) holds and coproducts are exact: if $f_i \colon A_i \to B_i$ is a
    family of monomorphisms, then $\bigoplus f_i \colon \bigoplus A_i \to \bigoplus B_i$
    is a monomorphism.
    \item[(AB4*)] (AB3*) holds and products are exact.
    \item[(AB5)] (AB3) holds and filtered colimits are exact: filtered colimits commute
    with finite limits.
    \item[(AB5*)] (AB3*) holds and filtered limits (cofiltered limits) are exact.
\end{enumerate}
A \textbf{Grothendieck category}\index{Grothendieck category} is an (AB5) abelian category
with a generator.
\end{definition}

\begin{example}\label{ex:grothendieck-cat}
\begin{enumerate}
    \item $\Mod_R$ satisfies (AB3)--(AB5) and (AB3*)--(AB4*). It is a Grothendieck category
    with generator $R$. It does \emph{not} satisfy (AB5*) in general.
    \item $\operatorname{Sh}(X, \Ab)$ is a Grothendieck category.
    \item The category $\Ab^{\mathrm{f}}$ of finite abelian groups satisfies (AB1)--(AB2)
    but not (AB3): infinite coproducts do not exist.
\end{enumerate}
\end{example}

\begin{theorem}[Freyd--Mitchell embedding]\label{thm:freyd-mitchell}
\index{Freyd--Mitchell embedding theorem}\index{embedding theorem}
Let $\mathcal{A}$ be a small abelian category. Then there exists a ring $R$ and an
exact, fully faithful functor $F \colon \mathcal{A} \hookrightarrow \Mod_R$.
\end{theorem}

\begin{remark}\label{rem:freyd-mitchell}
The Freyd--Mitchell theorem justifies ``diagram chasing'' in abstract abelian categories:
any statement about finite diagrams involving kernels, cokernels, exactness, etc., that
holds in $\Mod_R$ for all rings $R$ automatically holds in any abelian category.
We shall freely use element-wise arguments in proofs below, understood via this embedding.
\end{remark}


\section{Exact Sequences}\label{sec:exact-seq}

\begin{definition}[Exact sequence]\label{def:exact-seq}
\index{exact sequence}\index{sequence!exact}
Let $\mathcal{A}$ be an abelian category. A sequence of morphisms
\[
    \cdots \longrightarrow A_{n+1} \xrightarrow{f_{n+1}} A_n \xrightarrow{f_n} A_{n-1}
    \longrightarrow \cdots
\]
is \textbf{exact at $A_n$} if $\im(f_{n+1}) = \Ker(f_n)$ (as subobjects of $A_n$).
Equivalently, $f_n \circ f_{n+1} = 0$ and the induced morphism
$\operatorname{coim}(f_{n+1}) \to \Ker(f_n)$ is an isomorphism.
The sequence is \textbf{exact} if it is exact at every term.
\end{definition}

\begin{definition}[Short exact sequence]\label{def:ses}
\index{short exact sequence}\index{exact sequence!short}
A \textbf{short exact sequence} (SES) in $\mathcal{A}$ is an exact sequence
\[
    0 \longrightarrow A \xrightarrow{\;\;f\;\;} B \xrightarrow{\;\;g\;\;} C \longrightarrow 0.
\]
This means:
\begin{itemize}
    \item $f$ is a monomorphism (exactness at $A$),
    \item $g$ is an epimorphism (exactness at $C$),
    \item $\im(f) = \Ker(g)$ (exactness at $B$).
\end{itemize}
\end{definition}

\begin{example}\label{ex:ses-examples}
\index{short exact sequence!examples}
\begin{enumerate}
    \item In $\Ab$: $0 \to \mathbb{Z} \xrightarrow{\times n} \mathbb{Z} \to \mathbb{Z}/n\mathbb{Z} \to 0$
    for any $n \geq 1$.
    \item In $\Mod_R$: for any submodule $N \subseteq M$, we have
    $0 \to N \hookrightarrow M \twoheadrightarrow M/N \to 0$.
    \item A SES $0 \to A \to B \to C \to 0$ \textbf{splits}\index{split exact sequence}
    if there exists $s \colon C \to B$ with $g \circ s = \id_C$, or equivalently
    $r \colon B \to A$ with $r \circ f = \id_A$. In this case $B \cong A \oplus C$.
\end{enumerate}
\end{example}

\begin{proposition}[Splitting lemma]\label{prop:splitting-lemma}
\index{splitting lemma}
For a short exact sequence $0 \to A \xrightarrow{f} B \xrightarrow{g} C \to 0$
in an abelian category, the following are equivalent:
\begin{enumerate}
    \item There exists $s \colon C \to B$ with $g \circ s = \id_C$ (right splitting).
    \item There exists $r \colon B \to A$ with $r \circ f = \id_A$ (left splitting).
    \item $B \cong A \oplus C$ via an isomorphism compatible with $f$ and $g$.
\end{enumerate}
\end{proposition}

\begin{proof}
$(3) \Rightarrow (1)$ and $(3) \Rightarrow (2)$ are clear from the biproduct structure.

$(1) \Rightarrow (3)$: Given $s \colon C \to B$ with $g \circ s = \id_C$, define
$\varphi \colon A \oplus C \to B$ by $\varphi = f \circ \pi_A + s \circ \pi_C$.
Then $g \circ \varphi \circ \iota_C = g \circ s = \id_C$ and
$g \circ \varphi \circ \iota_A = g \circ f = 0$. We claim $\varphi$ is an isomorphism.
By the five lemma (Theorem~\ref{thm:five-lemma} below), or by direct argument:
for $b \in B$, write $b = f(a) + s(g(b))$ where $a$ is the unique element with
$f(a) = b - s(g(b))$ (which lies in $\Ker(g) = \im(f)$).

$(2) \Rightarrow (3)$: Dual argument.
\end{proof}


\section{Exact Functors}\label{sec:exact-functors}

\begin{definition}[Exact functor]\label{def:exact-functor}
\index{functor!exact}\index{exact functor}\index{functor!left exact}\index{functor!right exact}
Let $F \colon \mathcal{A} \to \mathcal{B}$ be an additive functor between abelian categories.
\begin{enumerate}
    \item $F$ is \textbf{left exact} if for every short exact sequence
    $0 \to A \to B \to C \to 0$, the sequence $0 \to FA \to FB \to FC$ is exact.
    \item $F$ is \textbf{right exact} if for every SES, $FA \to FB \to FC \to 0$ is exact.
    \item $F$ is \textbf{exact} if it is both left and right exact, i.e., it preserves
    short exact sequences.
\end{enumerate}
\end{definition}

\begin{example}\label{ex:exact-functors}
\begin{enumerate}
    \item\index{Hom@$\Hom$!left exactness}
    For any $R$-module $M$, the functor $\Hom_R(M, -)$ is left exact (covariant).
    The functor $\Hom_R(-, M)$ is left exact (contravariant, i.e., exact as a functor
    $\Mod_R^{\op} \to \Ab$).
    \item\index{tensor product!right exactness}
    The tensor product functor $M \otimes_R -$ is right exact for any right $R$-module $M$.
    \item The forgetful functor $\Mod_R \to \Ab$ is exact.
    \item\index{flat module}
    An $R$-module $M$ is \textbf{flat} if and only if $M \otimes_R -$ is exact.
\end{enumerate}
\end{example}

\begin{proposition}\label{prop:exact-functor-les}
\index{exact functor!characterizations}
Let $F \colon \mathcal{A} \to \mathcal{B}$ be an additive functor. Then:
\begin{enumerate}
    \item $F$ is left exact if and only if $F$ preserves kernels.
    \item $F$ is right exact if and only if $F$ preserves cokernels.
    \item $F$ is exact if and only if $F$ preserves short exact sequences.
\end{enumerate}
\end{proposition}


\section{The Snake Lemma}\label{sec:snake-lemma}

The snake lemma is one of the fundamental tools of homological algebra. It
produces long exact sequences from short ones.

\begin{theorem}[Snake lemma]\label{thm:snake-lemma}
\index{snake lemma}\index{connecting morphism}
Consider a commutative diagram in an abelian category $\mathcal{A}$ with exact rows:
\[
\begin{tikzcd}
    & A' \arrow[r, "f'"] \arrow[d, "a"] & B' \arrow[r, "g'"] \arrow[d, "b"] & C' \arrow[r] \arrow[d, "c"] & 0 \\
    0 \arrow[r] & A \arrow[r, "f"'] & B \arrow[r, "g"'] & C &
\end{tikzcd}
\]
Then there exists an exact sequence
\[
    \Ker(a) \xrightarrow{\;\Ker(f')\;} \Ker(b) \xrightarrow{\;\Ker(g')\;} \Ker(c)
    \xrightarrow{\;\;\delta\;\;} \coker(a) \xrightarrow{\;\coker(f)\;} \coker(b)
    \xrightarrow{\;\coker(g)\;} \coker(c),
\]
where $\delta$ is the \textbf{connecting morphism}\index{connecting morphism!snake lemma}.
Moreover:
\begin{itemize}
    \item If $f'$ is a monomorphism, so is $\Ker(a) \to \Ker(b)$.
    \item If $g$ is an epimorphism, so is $\coker(b) \to \coker(c)$.
\end{itemize}
\end{theorem}

\begin{proof}
We give a complete proof using element-chasing (justified by the Freyd--Mitchell theorem).

\medskip\noindent
\textbf{Step 1: The induced morphisms.}
The morphisms $\Ker(a) \to \Ker(b)$ and $\Ker(b) \to \Ker(c)$ are induced by $f'$ and $g'$
respectively. Indeed, if $x \in \Ker(a)$, then $b(f'(x)) = f(a(x)) = f(0) = 0$, so
$f'(x) \in \Ker(b)$. Similarly for $g'$. The morphisms $\coker(a) \to \coker(b) \to \coker(c)$
are induced dually.

\medskip\noindent
\textbf{Step 2: Construction of $\delta$.}
Let $x \in \Ker(c) \subseteq C'$. Since $g'$ is surjective (onto $C'$), choose
$y \in B'$ with $g'(y) = x$. Consider $b(y) \in B$. Then
\[
    g(b(y)) = c(g'(y)) = c(x) = 0,
\]
so $b(y) \in \Ker(g) = \im(f)$. Choose $z \in A$ with $f(z) = b(y)$ (unique since $f$
is injective). Define
\[
    \delta(x) = [z] \in \coker(a) = A / \im(a).
\]

\medskip\noindent
\textbf{Step 3: Well-definedness of $\delta$.}
If $y'$ is another choice with $g'(y') = x$, then $g'(y - y') = 0$, so $y - y' = f'(w)$
for some $w \in A'$. Then $b(y) - b(y') = b(f'(w)) = f(a(w))$, so $z - z' = a(w)$
since $f$ is injective. Thus $[z] = [z']$ in $\coker(a)$.

\medskip\noindent
\textbf{Step 4: Exactness at $\Ker(b)$.}
The composition $\Ker(a) \to \Ker(b) \to \Ker(c)$ is $g' \circ f' = 0$ (since the top
row is a complex). Conversely, let $y \in \Ker(b)$ with $g'(y) \in \Ker(c)$ mapping to $0$
in $\Ker(c)$---that is, $g'(y) = 0$. Then $y \in \Ker(g') = \im(f')$, so $y = f'(x)$
for some $x \in A'$. Now $f(a(x)) = b(f'(x)) = b(y) = 0$, and since $f$ is injective,
$a(x) = 0$, so $x \in \Ker(a)$.

\medskip\noindent
\textbf{Step 5: Exactness at $\Ker(c)$.}
The composition $\Ker(b) \to \Ker(c) \xrightarrow{\delta} \coker(a)$ is zero: if $y \in \Ker(b)$
and $x = g'(y) \in \Ker(c)$, then in the construction of $\delta$, we may choose $y$
itself, and $b(y) = 0$, so $z = 0$, giving $\delta(x) = 0$.

Conversely, let $x \in \Ker(c)$ with $\delta(x) = 0$. With notation as in Step~2,
$[z] = 0$ means $z = a(w)$ for some $w \in A'$. Then $b(y) = f(z) = f(a(w)) = b(f'(w))$,
so $y - f'(w) \in \Ker(b)$ and $g'(y - f'(w)) = g'(y) - 0 = x$.

\medskip\noindent
\textbf{Step 6: Exactness at $\coker(a)$.}
The composition $\delta$ followed by $\coker(a) \to \coker(b)$ is zero: in the
notation above, $\delta(x) = [z]$ and the image in $\coker(b)$ is $[f(z)] = [b(y)] = 0$.

Conversely, let $[z] \in \coker(a)$ with $[f(z)] = 0$ in $\coker(b)$. Then
$f(z) = b(y)$ for some $y \in B'$. Now $c(g'(y)) = g(b(y)) = g(f(z)) = 0$, so
$g'(y) \in \Ker(c)$, and $\delta(g'(y)) = [z]$ by construction.

\medskip\noindent
\textbf{Step 7: Exactness at $\coker(b)$.}
Similar argument, left as an exercise.

\medskip\noindent
\textbf{Step 8: Additional exactness.}
If $f'$ is mono, then $\Ker(a) \to \Ker(b)$ (given by $f'|_{\Ker(a)}$) is mono.
If $g$ is epi, then $\coker(b) \to \coker(c)$ is epi by a dual argument.
\end{proof}

The snake lemma is often remembered via the following diagram, where the ``snake''
traces the path through the kernels and cokernels:

\[
\begin{tikzcd}[column sep=1.2em, row sep=2.5em]
    & \Ker(a) \arrow[r] \arrow[d, hookrightarrow] & \Ker(b) \arrow[r] \arrow[d, hookrightarrow]
    & \Ker(c) \arrow[d, hookrightarrow] \arrow[dll, "\delta"', dashed,
    rounded corners, to path={
        -- ([xshift=2em]\tikztostart.east)
        -- ([xshift=2em, yshift=-2em]\tikztostart.east)
        -- ([xshift=-2em, yshift=-2em]\tikztotarget.west)
        -- ([xshift=-2em]\tikztotarget.west)
        -- (\tikztotarget.west)}] \\
    & A' \arrow[r, "f'"] \arrow[d, "a"] & B' \arrow[r, "g'"] \arrow[d, "b"]
    & C' \arrow[r] \arrow[d, "c"] & 0 \\
    0 \arrow[r] & A \arrow[r, "f"'] \arrow[d, twoheadrightarrow]
    & B \arrow[r, "g"'] \arrow[d, twoheadrightarrow] & C \arrow[d, twoheadrightarrow] \\
    & \coker(a) \arrow[r] & \coker(b) \arrow[r] & \coker(c)
\end{tikzcd}
\]


\section{The Five Lemma}\label{sec:five-lemma}

\begin{theorem}[Five lemma]\label{thm:five-lemma}
\index{five lemma}
Consider a commutative diagram with exact rows in an abelian category:
\[
\begin{tikzcd}
    A_1 \arrow[r, "f_1"] \arrow[d, "\alpha_1"'] & A_2 \arrow[r, "f_2"] \arrow[d, "\alpha_2"']
    & A_3 \arrow[r, "f_3"] \arrow[d, "\alpha_3"'] & A_4 \arrow[r, "f_4"] \arrow[d, "\alpha_4"']
    & A_5 \arrow[d, "\alpha_5"'] \\
    B_1 \arrow[r, "g_1"'] & B_2 \arrow[r, "g_2"'] & B_3 \arrow[r, "g_3"']
    & B_4 \arrow[r, "g_4"'] & B_5
\end{tikzcd}
\]
\begin{enumerate}
    \item If $\alpha_1$ is an epimorphism and $\alpha_2, \alpha_4$ are monomorphisms,
    then $\alpha_3$ is a monomorphism.
    \item If $\alpha_5$ is a monomorphism and $\alpha_2, \alpha_4$ are epimorphisms,
    then $\alpha_3$ is an epimorphism.
    \item If $\alpha_1$ is an epimorphism, $\alpha_5$ is a monomorphism, and
    $\alpha_2, \alpha_4$ are isomorphisms, then $\alpha_3$ is an isomorphism.
\end{enumerate}
\end{theorem}

\begin{proof}
We use element chasing (justified by the Freyd--Mitchell embedding).

\medskip\noindent
\textbf{Part (1):} Suppose $\alpha_3(a_3) = 0$ for some $a_3 \in A_3$. Then
$\alpha_4(f_3(a_3)) = g_3(\alpha_3(a_3)) = 0$. Since $\alpha_4$ is mono,
$f_3(a_3) = 0$. By exactness at $A_3$, there exists $a_2 \in A_2$ with $f_2(a_2) = a_3$.
Now $g_2(\alpha_2(a_2)) = \alpha_3(f_2(a_2)) = \alpha_3(a_3) = 0$. By exactness at $B_2$,
there exists $b_1 \in B_1$ with $g_1(b_1) = \alpha_2(a_2)$. Since $\alpha_1$ is epi,
$b_1 = \alpha_1(a_1)$ for some $a_1 \in A_1$. Then
\[
    \alpha_2(f_1(a_1)) = g_1(\alpha_1(a_1)) = g_1(b_1) = \alpha_2(a_2).
\]
Since $\alpha_2$ is mono, $f_1(a_1) = a_2$, so
$a_3 = f_2(a_2) = f_2(f_1(a_1)) = 0$ by exactness at $A_2$ (since $\im(f_1) = \Ker(f_2)$).
Thus $\alpha_3$ is mono.

\medskip\noindent
\textbf{Part (2):} Let $b_3 \in B_3$. Consider $g_3(b_3) \in B_4$. Since $\alpha_4$
is epi, there exists $a_4 \in A_4$ with $\alpha_4(a_4) = g_3(b_3)$. Now
$\alpha_5(f_4(a_4)) = g_4(\alpha_4(a_4)) = g_4(g_3(b_3)) = 0$ by exactness at $B_4$.
Since $\alpha_5$ is mono, $f_4(a_4) = 0$.
By exactness at $A_4$, there exists $a_3 \in A_3$ with $f_3(a_3) = a_4$.
Consider $b_3 - \alpha_3(a_3)$:
\[
    g_3(b_3 - \alpha_3(a_3)) = g_3(b_3) - g_3(\alpha_3(a_3))
    = g_3(b_3) - \alpha_4(f_3(a_3)) = g_3(b_3) - \alpha_4(a_4) = 0.
\]
By exactness at $B_3$, there exists $b_2 \in B_2$ with $g_2(b_2) = b_3 - \alpha_3(a_3)$.
Since $\alpha_2$ is epi, $b_2 = \alpha_2(a_2)$ for some $a_2 \in A_2$. Then
\[
    b_3 = \alpha_3(a_3) + g_2(\alpha_2(a_2)) = \alpha_3(a_3) + \alpha_3(f_2(a_2))
    = \alpha_3(a_3 + f_2(a_2)).
\]
Thus $\alpha_3$ is epi.

\medskip\noindent
\textbf{Part (3):} Immediate from (1) and (2).
\end{proof}

\begin{corollary}[Short five lemma]\label{cor:short-five-lemma}
\index{five lemma!short}
Given a commutative diagram with exact rows
\[
\begin{tikzcd}
    0 \arrow[r] & A' \arrow[r] \arrow[d, "\alpha"'] & B' \arrow[r] \arrow[d, "\beta"'] & C' \arrow[r] \arrow[d, "\gamma"'] & 0 \\
    0 \arrow[r] & A \arrow[r] & B \arrow[r] & C \arrow[r] & 0
\end{tikzcd}
\]
if $\alpha$ and $\gamma$ are isomorphisms, then $\beta$ is an isomorphism.
\end{corollary}

\begin{proof}
Apply the five lemma with $A_1 = B_1 = 0$ and $A_5 = B_5 = 0$, taking
$\alpha_1 = \alpha_5 = 0$, $\alpha_2 = \alpha$, $\alpha_3 = \beta$, $\alpha_4 = \gamma$.
\end{proof}


\section{Injective and Projective Objects}\label{sec:inj-proj}

\begin{definition}[Projective object]\label{def:projective}
\index{projective object}\index{object!projective}
An object $P$ in an abelian category $\mathcal{A}$ is \textbf{projective} if the
functor $\Hom_{\mathcal{A}}(P, -)$ is exact. Equivalently, for every epimorphism
$p \colon B \twoheadrightarrow C$ and every morphism $f \colon P \to C$, there exists
a morphism $\tilde{f} \colon P \to B$ with $p \circ \tilde{f} = f$:
\[
\begin{tikzcd}
    & P \arrow[d, "f"] \arrow[dl, dashed, "\tilde{f}"'] \\
    B \arrow[r, "p"', twoheadrightarrow] & C \arrow[r] & 0
\end{tikzcd}
\]
\end{definition}

\begin{definition}[Injective object]\label{def:injective}
\index{injective object}\index{object!injective}
An object $I$ in $\mathcal{A}$ is \textbf{injective} if $\Hom_{\mathcal{A}}(-, I)$ is exact.
Equivalently, for every monomorphism $i \colon A \hookrightarrow B$ and every
$f \colon A \to I$, there exists $\tilde{f} \colon B \to I$ with $\tilde{f} \circ i = f$:
\[
\begin{tikzcd}
    0 \arrow[r] & A \arrow[r, hookrightarrow, "i"] \arrow[d, "f"'] & B \arrow[dl, dashed, "\tilde{f}"] \\
    & I &
\end{tikzcd}
\]
\end{definition}

\begin{example}\label{ex:proj-inj-examples}
\begin{enumerate}
    \item\index{projective module}
    In $\Mod_R$, free modules are projective. A module is projective if and only if it
    is a direct summand of a free module.
    \item\index{injective!abelian group}
    In $\Ab$, the group $\mathbb{Q}$ is injective. More generally, $\mathbb{Q}/\mathbb{Z}$
    is injective.
    \item In a semisimple abelian category (every SES splits), every object is both
    projective and injective.
\end{enumerate}
\end{example}

\begin{definition}[Enough projectives/injectives]\label{def:enough-proj-inj}
\index{enough projectives}\index{enough injectives}
An abelian category $\mathcal{A}$ has \textbf{enough projectives} if for every object $A$,
there exists an epimorphism $P \twoheadrightarrow A$ with $P$ projective. It has
\textbf{enough injectives} if for every $A$, there exists a monomorphism $A \hookrightarrow I$
with $I$ injective.
\end{definition}

\begin{theorem}[Baer's criterion]\label{thm:baer-criterion}
\index{Baer's criterion}
Let $R$ be a ring and $I$ an $R$-module. Then $I$ is injective if and only if for every
left ideal $J \subseteq R$ and every homomorphism $f \colon J \to I$, there exists an
extension $\tilde{f} \colon R \to I$ with $\tilde{f}|_J = f$.
\end{theorem}

\begin{theorem}\label{thm:grothendieck-injectives}
\index{Grothendieck category!enough injectives}
Every Grothendieck category has enough injectives. In particular, $\Mod_R$ and
$\operatorname{Sh}(X, \Ab)$ have enough injectives.
\end{theorem}


\section{Resolutions}\label{sec:resolutions}

\begin{definition}[Projective resolution]\label{def:proj-resolution}
\index{projective resolution}\index{resolution!projective}
A \textbf{projective resolution} of an object $A$ in an abelian category $\mathcal{A}$
is an exact sequence
\[
    \cdots \longrightarrow P_2 \xrightarrow{d_2} P_1 \xrightarrow{d_1} P_0
    \xrightarrow{\varepsilon} A \longrightarrow 0
\]
where each $P_n$ is projective. We write this as $P_\bullet \xrightarrow{\varepsilon} A$.
\end{definition}

\begin{definition}[Injective resolution]\label{def:inj-resolution}
\index{injective resolution}\index{resolution!injective}
An \textbf{injective resolution} of $A$ is an exact sequence
\[
    0 \longrightarrow A \xrightarrow{\eta} I^0 \xrightarrow{d^0} I^1
    \xrightarrow{d^1} I^2 \longrightarrow \cdots
\]
where each $I^n$ is injective. We write this as $A \xrightarrow{\eta} I^\bullet$.
\end{definition}

\begin{proposition}\label{prop:resolution-existence}
If $\mathcal{A}$ has enough projectives, then every object admits a projective resolution.
If $\mathcal{A}$ has enough injectives, then every object admits an injective resolution.
\end{proposition}

\begin{proof}
We construct a projective resolution inductively. Choose an epi $P_0 \twoheadrightarrow A$
with $P_0$ projective and let $K_0 = \Ker(P_0 \to A)$. Then choose $P_1 \twoheadrightarrow K_0$
with $P_1$ projective, set $d_1 \colon P_1 \twoheadrightarrow K_0 \hookrightarrow P_0$,
let $K_1 = \Ker(P_1 \to K_0)$, and continue.
\end{proof}

\begin{theorem}[Comparison theorem]\label{thm:comparison}
\index{comparison theorem!for resolutions}
Let $P_\bullet \xrightarrow{\varepsilon} A$ be a projective resolution and
$Q_\bullet \xrightarrow{\varepsilon'} B$ be any resolution (not necessarily projective).
Then any morphism $f \colon A \to B$ lifts to a chain map
$f_\bullet \colon P_\bullet \to Q_\bullet$:
\[
\begin{tikzcd}
    \cdots \arrow[r] & P_2 \arrow[r, "d_2"] \arrow[d, dashed, "f_2"'] & P_1 \arrow[r, "d_1"] \arrow[d, dashed, "f_1"'] & P_0 \arrow[r, "\varepsilon"] \arrow[d, dashed, "f_0"'] & A \arrow[r] \arrow[d, "f"] & 0 \\
    \cdots \arrow[r] & Q_2 \arrow[r, "d_2'"'] & Q_1 \arrow[r, "d_1'"'] & Q_0 \arrow[r, "\varepsilon'"'] & B \arrow[r] & 0
\end{tikzcd}
\]
Moreover, any two such lifts are chain homotopic.
\end{theorem}

\begin{proof}
\textit{Existence.} We construct $f_n$ inductively. The composite $f \circ \varepsilon \colon P_0 \to B$
factors through $\varepsilon' \colon Q_0 \to B$ since $P_0$ is projective and $\varepsilon'$
is epi. This gives $f_0$. For the inductive step, given $f_{n-1}$, we need $f_n$ making the
square commute. The morphism $f_{n-1} \circ d_n \colon P_n \to Q_{n-1}$ satisfies
$d'_{n-1} \circ f_{n-1} \circ d_n = f_{n-2} \circ d_{n-1} \circ d_n = 0$. Thus
$f_{n-1} \circ d_n$ factors through $\Ker(d'_{n-1}) = \im(d'_n)$, and projectivity of $P_n$
gives $f_n$.

\textit{Uniqueness up to homotopy.} Suppose $f_\bullet, g_\bullet$ are two lifts.
We construct a chain homotopy $s_n \colon P_n \to Q_{n+1}$ with
$f_n - g_n = d'_{n+1} s_n + s_{n-1} d_n$ inductively, using projectivity at each step.
\end{proof}


\section{Exercises}\label{sec:abelian-exercises}

\begin{exercise}\label{exo:preadditive-zero}
\index{zero object!in preadditive category}
Let $\mathcal{A}$ be a preadditive category. Show that $\mathcal{A}$ has a zero object if and
only if it has an object $Z$ such that $\id_Z = 0_{Z,Z}$.
\end{exercise}

\begin{exercise}\label{exo:additive-functor-zero}
Show that an additive functor preserves zero objects, biproducts, kernels, and cokernels
(when they exist).
\end{exercise}

\begin{exercise}\label{exo:idempotent-split}
\index{idempotent}\index{Karoubi envelope}
Let $\mathcal{A}$ be an additive category and $e \colon A \to A$ an idempotent ($e^2 = e$).
We say $e$ \textbf{splits} if there exist morphisms $p \colon A \to B$ and $i \colon B \to A$
with $p \circ i = \id_B$ and $i \circ p = e$.
\begin{enumerate}
    \item Show that in an abelian category, every idempotent splits.
    \item Show that $B = \Ker(\id_A - e)$ works.
\end{enumerate}
\end{exercise}

\begin{exercise}\label{exo:nine-lemma}
\index{nine lemma}\index{$3 \times 3$ lemma}
(The $3 \times 3$ lemma.) Consider a commutative diagram with exact columns and
the middle and bottom rows exact:
\[
\begin{tikzcd}
    & 0 \arrow[d] & 0 \arrow[d] & 0 \arrow[d] & \\
    0 \arrow[r] & A' \arrow[r] \arrow[d] & B' \arrow[r] \arrow[d] & C' \arrow[r] \arrow[d] & 0 \\
    0 \arrow[r] & A \arrow[r] \arrow[d] & B \arrow[r] \arrow[d] & C \arrow[r] \arrow[d] & 0 \\
    0 \arrow[r] & A'' \arrow[r] \arrow[d] & B'' \arrow[r] \arrow[d] & C'' \arrow[r] \arrow[d] & 0 \\
    & 0 & 0 & 0 &
\end{tikzcd}
\]
Prove that the top row is also exact.
\end{exercise}

\begin{exercise}\label{exo:horseshoe}
\index{horseshoe lemma}
(Horseshoe lemma.) Let $0 \to A' \to A \to A'' \to 0$ be a short exact sequence in an
abelian category with enough projectives. Given projective resolutions
$P'_\bullet \to A'$ and $P''_\bullet \to A''$, show there exists a projective resolution
$P_\bullet \to A$ with $P_n = P'_n \oplus P''_n$ fitting into a short exact sequence
of chain complexes $0 \to P'_\bullet \to P_\bullet \to P''_\bullet \to 0$.
\end{exercise}

\begin{exercise}\label{exo:injective-divisible}
\index{divisible group}\index{injective!abelian group}
Show that an abelian group $G$ is injective in $\Ab$ if and only if it is divisible
(for every $g \in G$ and $n \geq 1$, there exists $h \in G$ with $nh = g$).
\end{exercise}

\begin{exercise}\label{exo:projective-characterization}
\index{projective module!characterization}
Let $R$ be a ring and $P$ an $R$-module. Prove the equivalence:
\begin{enumerate}
    \item $P$ is projective.
    \item Every short exact sequence $0 \to A \to B \to P \to 0$ splits.
    \item $P$ is a direct summand of a free module.
\end{enumerate}
\end{exercise}

\begin{exercise}\label{exo:abelian-pullback}
Let $\mathcal{A}$ be an abelian category. Show that the pullback of an epimorphism is
an epimorphism, and the pushout of a monomorphism is a monomorphism.
\end{exercise}


%% ========================================================================
\chapter{Derived Categories --- Introduction}\label{ch:derived}
\minitoc

\section{Chain Complexes and Their Morphisms}\label{sec:chain-complexes}

\begin{definition}[Chain complex]\label{def:chain-complex}
\index{chain complex}\index{complex!chain}
Let $\mathcal{A}$ be an abelian category. A \textbf{chain complex} (or simply \textbf{complex})
in $\mathcal{A}$ is a sequence of objects and morphisms
\[
    A^\bullet \colon \quad \cdots \longrightarrow A^{n-1} \xrightarrow{d^{n-1}} A^n
    \xrightarrow{d^n} A^{n+1} \longrightarrow \cdots
\]
such that $d^n \circ d^{n-1} = 0$ for all $n \in \mathbb{Z}$. The morphisms $d^n$ are
called \textbf{differentials}\index{differential}. We use cohomological (superscript)
indexing throughout.
\end{definition}

\begin{definition}[Cohomology]\label{def:cohomology-complex}
\index{cohomology!of a complex}
The $n$-th \textbf{cohomology} of a complex $A^\bullet$ is
\[
    H^n(A^\bullet) = \Ker(d^n) / \im(d^{n-1}) = \Ker(d^n \colon A^n \to A^{n+1}) /
    \im(d^{n-1} \colon A^{n-1} \to A^n).
\]
A complex is \textbf{acyclic}\index{acyclic complex} (or \textbf{exact}) if $H^n(A^\bullet) = 0$
for all $n$.
\end{definition}

\begin{definition}[Category of complexes]\label{def:complex-category}
\index{Ch(A)@$\operatorname{Ch}(\mathcal{A})$}
\index{category!of complexes}
The \textbf{category of chain complexes} $\operatorname{Ch}(\mathcal{A})$ has:
\begin{itemize}
    \item Objects: chain complexes in $\mathcal{A}$.
    \item Morphisms: a \textbf{chain map}\index{chain map} $f \colon A^\bullet \to B^\bullet$
    is a collection of morphisms $f^n \colon A^n \to B^n$ commuting with differentials:
    $d_B^n \circ f^n = f^{n+1} \circ d_A^n$ for all $n$.
\end{itemize}
We also define the bounded variants:
\begin{align*}
    \operatorname{Ch}^+(\mathcal{A}) &= \{ A^\bullet \mid A^n = 0 \text{ for } n \ll 0 \}
    &&\text{(bounded below)}, \\
    \operatorname{Ch}^-(\mathcal{A}) &= \{ A^\bullet \mid A^n = 0 \text{ for } n \gg 0 \}
    &&\text{(bounded above)}, \\
    \operatorname{Ch}^b(\mathcal{A}) &= \operatorname{Ch}^+(\mathcal{A}) \cap \operatorname{Ch}^-(\mathcal{A})
    &&\text{(bounded)}.
\end{align*}
\end{definition}

\begin{remark}\label{rem:ch-abelian}
\index{Ch(A)@$\operatorname{Ch}(\mathcal{A})$!abelian structure}
The category $\operatorname{Ch}(\mathcal{A})$ is abelian when $\mathcal{A}$ is abelian,
with kernels, cokernels, and exact sequences computed degreewise.
\end{remark}

\begin{definition}[Shift functor]\label{def:shift-functor}
\index{shift functor}\index{translation functor}
The \textbf{shift} (or \textbf{translation}) functor $[1] \colon \operatorname{Ch}(\mathcal{A})
\to \operatorname{Ch}(\mathcal{A})$ is defined by
\[
    (A[1])^n = A^{n+1}, \qquad d_{A[1]}^n = -d_A^{n+1}.
\]
More generally, $A[k]^n = A^{n+k}$ with $d_{A[k]}^n = (-1)^k d_A^{n+k}$.
For a chain map $f \colon A^\bullet \to B^\bullet$, we set $f[1]^n = f^{n+1}$.
\end{definition}

\begin{definition}[Mapping cone]\label{def:mapping-cone}
\index{mapping cone}\index{cone!of a morphism}
Let $f \colon A^\bullet \to B^\bullet$ be a chain map. The \textbf{mapping cone}
$\operatorname{Cone}(f)$ is the complex defined by
\[
    \operatorname{Cone}(f)^n = A^{n+1} \oplus B^n, \qquad
    d_{\operatorname{Cone}(f)}^n = \begin{pmatrix} -d_A^{n+1} & 0 \\ f^{n+1} & d_B^n \end{pmatrix}.
\]
One verifies $d^2 = 0$ using $d_B \circ f = f \circ d_A$.
There is a natural short exact sequence of complexes:
\[
    0 \longrightarrow B^\bullet \xrightarrow{\;\iota\;} \operatorname{Cone}(f)
    \xrightarrow{\;\pi\;} A[1]^\bullet \longrightarrow 0.
\]
\end{definition}

\begin{proposition}\label{prop:cone-les}
\index{mapping cone!long exact sequence}
For any chain map $f \colon A^\bullet \to B^\bullet$, the mapping cone gives a long exact
sequence in cohomology:
\[
    \cdots \to H^n(A^\bullet) \xrightarrow{H^n(f)} H^n(B^\bullet) \to
    H^n(\operatorname{Cone}(f)) \to H^{n+1}(A^\bullet) \xrightarrow{H^{n+1}(f)} \cdots
\]
In particular, $f$ is a quasi-isomorphism if and only if $\operatorname{Cone}(f)$ is acyclic.
\end{proposition}


\section{Chain Homotopies and the Homotopy Category}\label{sec:homotopy-cat}

\begin{definition}[Chain homotopy]\label{def:chain-homotopy}
\index{chain homotopy}\index{homotopy!of chain maps}
Let $f, g \colon A^\bullet \to B^\bullet$ be chain maps. A \textbf{chain homotopy}
from $f$ to $g$ is a collection of morphisms $s^n \colon A^n \to B^{n-1}$ such that
\[
    f^n - g^n = d_B^{n-1} \circ s^n + s^{n+1} \circ d_A^n \quad \text{for all } n.
\]
\[
\begin{tikzcd}[column sep=2em]
    \cdots \arrow[r] & A^{n-1} \arrow[r, "d^{n-1}"] \arrow[d, shift left=1, "f^{n-1}"]
    \arrow[d, shift right=1, "g^{n-1}"']
    & A^n \arrow[r, "d^n"] \arrow[d, shift left=1, "f^n"]
    \arrow[d, shift right=1, "g^n"'] \arrow[ld, dashed, "s^n"']
    & A^{n+1} \arrow[r] \arrow[d, shift left=1, "f^{n+1}"]
    \arrow[d, shift right=1, "g^{n+1}"'] \arrow[ld, dashed, "s^{n+1}"']
    & \cdots \\
    \cdots \arrow[r] & B^{n-1} \arrow[r, "d^{n-1}"'] & B^n \arrow[r, "d^n"'] & B^{n+1} \arrow[r] & \cdots
\end{tikzcd}
\]
We write $f \sim g$ and say $f$ and $g$ are \textbf{homotopic}.
If $f \sim 0$, we say $f$ is \textbf{null-homotopic}\index{null-homotopic}.
\end{definition}

\begin{proposition}\label{prop:homotopy-equiv-rel}
Chain homotopy is an equivalence relation on $\Hom_{\operatorname{Ch}(\mathcal{A})}(A^\bullet, B^\bullet)$,
compatible with composition: if $f \sim g$ then $h \circ f \sim h \circ g$ and
$f \circ k \sim g \circ k$ for any compatible chain maps $h, k$.
\end{proposition}

\begin{proof}
Reflexivity: take $s = 0$. Symmetry: negate $s$. Transitivity: add homotopies.
For compatibility, if $s$ is a homotopy from $f$ to $g$, then $h \circ s$ is a homotopy
from $h \circ f$ to $h \circ g$, and $s \circ k$ is a homotopy from $f \circ k$ to $g \circ k$.
\end{proof}

\begin{proposition}\label{prop:homotopic-cohomology}
\index{chain homotopy!and cohomology}
If $f \sim g \colon A^\bullet \to B^\bullet$, then $H^n(f) = H^n(g)$ for all $n$.
\end{proposition}

\begin{proof}
For $[a] \in H^n(A^\bullet)$ with $d_A^n(a) = 0$:
$f^n(a) - g^n(a) = d_B^{n-1}(s^n(a)) + s^{n+1}(d_A^n(a)) = d_B^{n-1}(s^n(a))$,
which is a coboundary. Thus $[f^n(a)] = [g^n(a)]$.
\end{proof}

\begin{definition}[Homotopy category]\label{def:homotopy-category}
\index{homotopy category}\index{K(A)@$K(\mathcal{A})$}
The \textbf{homotopy category} $K(\mathcal{A})$ has:
\begin{itemize}
    \item Objects: chain complexes in $\mathcal{A}$ (same as $\operatorname{Ch}(\mathcal{A})$).
    \item Morphisms: $\Hom_{K(\mathcal{A})}(A^\bullet, B^\bullet) =
    \Hom_{\operatorname{Ch}(\mathcal{A})}(A^\bullet, B^\bullet) / {\sim}$,
    i.e., chain maps modulo homotopy.
\end{itemize}
The bounded variants $K^+(\mathcal{A})$, $K^-(\mathcal{A})$, $K^b(\mathcal{A})$ are
defined similarly.
\end{definition}

\begin{remark}\label{rem:homotopy-not-abelian}
\index{homotopy category!not abelian}
The homotopy category $K(\mathcal{A})$ is additive but \emph{not} abelian in general.
It does, however, carry a \textbf{triangulated structure}\index{triangulated category}
(see Section~\ref{sec:triangulated}).
\end{remark}

\begin{definition}[Homotopy equivalence]\label{def:homotopy-equivalence}
\index{homotopy equivalence}
A chain map $f \colon A^\bullet \to B^\bullet$ is a \textbf{homotopy equivalence} if
it becomes an isomorphism in $K(\mathcal{A})$, i.e., there exists $g \colon B^\bullet \to A^\bullet$
with $g \circ f \sim \id_A$ and $f \circ g \sim \id_B$.
\end{definition}


\section{Quasi-Isomorphisms and the Derived Category}\label{sec:derived-cat}

\begin{definition}[Quasi-isomorphism]\label{def:quasi-iso}
\index{quasi-isomorphism}
A chain map $f \colon A^\bullet \to B^\bullet$ is a \textbf{quasi-isomorphism} if the induced
maps $H^n(f) \colon H^n(A^\bullet) \to H^n(B^\bullet)$ are isomorphisms for all $n \in \mathbb{Z}$.
We denote the class of quasi-isomorphisms by $\operatorname{qis}$.
\end{definition}

\begin{remark}\label{rem:homotopy-vs-quasi}
Every homotopy equivalence is a quasi-isomorphism (by Proposition~\ref{prop:homotopic-cohomology}),
but the converse is false in general. The derived category is obtained by formally
inverting all quasi-isomorphisms.
\end{remark}

\begin{example}\label{ex:quasi-iso-not-homotopy}
\index{quasi-isomorphism!vs.\ homotopy equivalence}
In $\Ab$, consider the complexes $A^\bullet \colon 0 \to \mathbb{Z} \xrightarrow{2} \mathbb{Z} \to 0$
(in degrees 0 and 1) and $B^\bullet \colon 0 \to \mathbb{Z}/2\mathbb{Z} \to 0$ (in degree 1).
The natural projection $A^\bullet \to B^\bullet$ is a quasi-isomorphism but not a homotopy
equivalence (since $\Hom(\mathbb{Z}/2\mathbb{Z}, \mathbb{Z}) = 0$, there is no map backward
in degree 1).
\end{example}

\subsection{Localization of categories}\label{subsec:localization}

\begin{definition}[Localization]\label{def:localization-cat}
\index{localization!of categories}\index{category!localization}
Let $\mathcal{C}$ be a category and $S$ a class of morphisms in $\mathcal{C}$.
The \textbf{localization} $\mathcal{C}[S^{-1}]$ is a category equipped with a functor
$Q \colon \mathcal{C} \to \mathcal{C}[S^{-1}]$ such that:
\begin{enumerate}
    \item $Q(s)$ is an isomorphism for every $s \in S$.
    \item $Q$ is universal with this property: for any functor $F \colon \mathcal{C} \to \mathcal{D}$
    sending $S$ to isomorphisms, there exists a unique $\bar{F} \colon \mathcal{C}[S^{-1}] \to \mathcal{D}$
    with $\bar{F} \circ Q = F$.
\end{enumerate}
\[
\begin{tikzcd}
    \mathcal{C} \arrow[r, "Q"] \arrow[rd, "F"'] & \mathcal{C}[S^{-1}] \arrow[d, dashed, "\exists!\bar{F}"] \\
    & \mathcal{D}
\end{tikzcd}
\]
\end{definition}

\begin{remark}\label{rem:localization-construction}
\index{localization!Gabriel--Zisman}
The localization always exists (Gabriel--Zisman): objects of $\mathcal{C}[S^{-1}]$ are those
of $\mathcal{C}$, and morphisms are equivalence classes of ``zigzag'' chains
\[
    A = C_0 \xleftarrow{s_1} C_1 \xrightarrow{f_1} C_2 \xleftarrow{s_2}
    C_3 \xrightarrow{f_2} \cdots \to B,
\]
where the backward arrows $s_i$ belong to $S$, subject to appropriate relations.
In general, $\Hom$ sets may fail to be sets (set-theoretic issues). However,
when $S$ admits a \textbf{calculus of fractions}\index{calculus of fractions},
morphisms simplify to ``roofs'' and the issues become manageable.
\end{remark}

\begin{definition}[Right/left roof]\label{def:roof}
\index{roof}\index{calculus of fractions}
When $S$ admits a \textbf{calculus of right fractions}, every morphism in $\mathcal{C}[S^{-1}]$
from $A$ to $B$ can be represented by a \textbf{right roof}:
\[
\begin{tikzcd}
    & C \arrow[dl, "s"'] \arrow[dr, "f"] & \\
    A & & B
\end{tikzcd}
\]
where $s \in S$. This represents the ``fraction'' $f \circ s^{-1}$.
Two roofs $(C, s, f)$ and $(C', s', f')$ represent the same morphism if and only if
they can be ``completed'' to a common refinement.
\end{definition}

\subsection{The derived category}\label{subsec:derived-cat-def}

\begin{definition}[Derived category]\label{def:derived-cat}
\index{derived category}\index{D(A)@$D(\mathcal{A})$}
Let $\mathcal{A}$ be an abelian category. The \textbf{derived category} $D(\mathcal{A})$
is the localization of $K(\mathcal{A})$ (the homotopy category) at the class of
quasi-isomorphisms:
\[
    D(\mathcal{A}) = K(\mathcal{A})[\operatorname{qis}^{-1}].
\]
The bounded variants are:
\begin{align*}
    D^+(\mathcal{A}) &= K^+(\mathcal{A})[\operatorname{qis}^{-1}], \\
    D^-(\mathcal{A}) &= K^-(\mathcal{A})[\operatorname{qis}^{-1}], \\
    D^b(\mathcal{A}) &= K^b(\mathcal{A})[\operatorname{qis}^{-1}].
\end{align*}
\end{definition}

\begin{remark}\label{rem:why-two-steps}
\index{derived category!construction}
The construction proceeds in two steps: $\operatorname{Ch}(\mathcal{A}) \to K(\mathcal{A}) \to D(\mathcal{A})$.
The first step (modding out homotopies) ensures that quasi-isomorphisms satisfy the Ore
conditions (calculus of fractions), making the second localization well-behaved.
If one tried to localize $\operatorname{Ch}(\mathcal{A})$ directly at quasi-isomorphisms,
the calculus of fractions would fail.
\end{remark}

\begin{theorem}\label{thm:derived-cat-roofs}
\index{derived category!morphisms as roofs}
The class of quasi-isomorphisms in $K(\mathcal{A})$ admits both a left and a right calculus
of fractions. Therefore, morphisms in $D(\mathcal{A})$ from $A^\bullet$ to $B^\bullet$ are
represented by roofs
\[
\begin{tikzcd}
    & C^\bullet \arrow[dl, "s"', "\sim"] \arrow[dr, "f"] & \\
    A^\bullet & & B^\bullet
\end{tikzcd}
\]
where $s$ is a quasi-isomorphism in $K(\mathcal{A})$ and $f$ is any morphism in $K(\mathcal{A})$.
\end{theorem}

\begin{proposition}\label{prop:derived-cat-embedding}
\index{derived category!embedding of $\mathcal{A}$}
There is a fully faithful functor $\mathcal{A} \hookrightarrow D(\mathcal{A})$ sending
an object $A$ to the complex concentrated in degree zero:
\[
    A \mapsto (\cdots \to 0 \to A \to 0 \to \cdots).
\]
Under this embedding, $\Hom_{D(\mathcal{A})}(A, B) \cong \Hom_{\mathcal{A}}(A, B)$
for objects $A, B \in \mathcal{A}$.
\end{proposition}

\begin{theorem}\label{thm:injective-resolution-equiv}
\index{derived category!via injective resolutions}
Let $\mathcal{A}$ be an abelian category with enough injectives, and let
$K^+(\mathcal{I})$ be the full subcategory of $K^+(\mathcal{A})$ consisting of
complexes of injective objects. Then the natural functor
\[
    K^+(\mathcal{I}) \xrightarrow{\;\sim\;} D^+(\mathcal{A})
\]
is an equivalence of categories. Dually, if $\mathcal{A}$ has enough projectives,
$K^-(\mathcal{P}) \xrightarrow{\sim} D^-(\mathcal{A})$.
\end{theorem}

\begin{proof}[Proof sketch]
The functor is the composition $K^+(\mathcal{I}) \hookrightarrow K^+(\mathcal{A})
\xrightarrow{Q} D^+(\mathcal{A})$.

\textit{Essential surjectivity:} Every bounded below complex $A^\bullet$ admits a
quasi-isomorphism $A^\bullet \xrightarrow{\sim} I^\bullet$ to a complex of injectives
(by inductively resolving, using an injective Cartan--Eilenberg resolution).

\textit{Fully faithfulness:} For complexes of injectives $I^\bullet$, $J^\bullet$,
a quasi-isomorphism $s \colon I^\bullet \xrightarrow{\sim} J^\bullet$ is already a
homotopy equivalence (a standard argument using injectivity at each degree).
Thus roofs reduce to ordinary morphisms in $K^+(\mathcal{I})$.
\end{proof}


\section{Triangulated Structure}\label{sec:triangulated}

\begin{definition}[Triangulated category]\label{def:triangulated-cat}
\index{triangulated category}\index{distinguished triangle}
A \textbf{triangulated category} is an additive category $\mathcal{T}$ equipped with:
\begin{enumerate}
    \item An automorphism $[1] \colon \mathcal{T} \to \mathcal{T}$ (the \textbf{shift functor}).
    \item A class of \textbf{distinguished triangles} (or \textbf{exact triangles})
    \[
        A \xrightarrow{f} B \xrightarrow{g} C \xrightarrow{h} A[1],
    \]
\end{enumerate}
satisfying axioms (TR1)--(TR4):
\begin{enumerate}
    \item[(TR1)] (a) For every $A$, the triangle $A \xrightarrow{\id} A \to 0 \to A[1]$ is
    distinguished. (b) Every morphism $f \colon A \to B$ can be completed to a distinguished
    triangle $A \xrightarrow{f} B \to C \to A[1]$. (c) A triangle isomorphic to a
    distinguished triangle is distinguished.
    \item[(TR2)] (\textbf{Rotation})\index{rotation of triangles}
    $A \xrightarrow{f} B \xrightarrow{g} C \xrightarrow{h} A[1]$ is distinguished
    if and only if $B \xrightarrow{g} C \xrightarrow{h} A[1] \xrightarrow{-f[1]} B[1]$ is.
    \item[(TR3)] (\textbf{Morphism of triangles}) Given distinguished triangles and morphisms
    $\alpha, \beta$ making the left square commute, there exists (not necessarily unique)
    $\gamma$ completing the morphism of triangles:
    \[
    \begin{tikzcd}
        A \arrow[r, "f"] \arrow[d, "\alpha"'] & B \arrow[r, "g"] \arrow[d, "\beta"']
        & C \arrow[r, "h"] \arrow[d, dashed, "\gamma"'] & A{[1]} \arrow[d, "\alpha{[1]}"'] \\
        A' \arrow[r, "f'"'] & B' \arrow[r, "g'"'] & C' \arrow[r, "h'"'] & A'{[1]}
    \end{tikzcd}
    \]
    \item[(TR4)] (\textbf{Octahedron axiom})\index{octahedron axiom}
    Given composable morphisms $A \xrightarrow{f} B \xrightarrow{g} C$ and distinguished
    triangles completing $f$, $g$, and $g \circ f$, there exists a distinguished triangle
    relating the three cones, making all faces commute.
\end{enumerate}
\end{definition}

\begin{theorem}\label{thm:K-A-triangulated}
\index{homotopy category!triangulated structure}
The homotopy category $K(\mathcal{A})$ is triangulated with:
\begin{itemize}
    \item Shift: $A^\bullet \mapsto A[1]^\bullet$ (Definition~\ref{def:shift-functor}).
    \item Distinguished triangles: those isomorphic in $K(\mathcal{A})$ to
    \[
        A^\bullet \xrightarrow{f} B^\bullet \xrightarrow{\iota} \operatorname{Cone}(f)
        \xrightarrow{\pi} A[1]^\bullet
    \]
    for some chain map $f$.
\end{itemize}
\end{theorem}

\begin{theorem}\label{thm:D-A-triangulated}
\index{derived category!triangulated structure}
The derived category $D(\mathcal{A})$ inherits a triangulated structure from $K(\mathcal{A})$
via the localization functor $Q \colon K(\mathcal{A}) \to D(\mathcal{A})$: a triangle in
$D(\mathcal{A})$ is distinguished if it is isomorphic to the image of a distinguished
triangle in $K(\mathcal{A})$.
\end{theorem}

\begin{proposition}\label{prop:ses-to-triangle}
\index{short exact sequence!and distinguished triangles}
A short exact sequence of complexes $0 \to A^\bullet \xrightarrow{f} B^\bullet
\xrightarrow{g} C^\bullet \to 0$ gives rise to a distinguished triangle
\[
    A^\bullet \xrightarrow{f} B^\bullet \xrightarrow{g} C^\bullet
    \xrightarrow{\delta} A[1]^\bullet
\]
in $D(\mathcal{A})$, where $\delta$ is the connecting morphism. In particular, the
long exact sequence in cohomology associated to a short exact sequence of complexes
is a consequence of the triangulated structure.
\end{proposition}

\begin{proposition}\label{prop:cohomological-functor}
\index{cohomological functor}
Let $\mathcal{T}$ be a triangulated category and $\mathcal{A}$ an abelian category.
An additive functor $H \colon \mathcal{T} \to \mathcal{A}$ is \textbf{cohomological} if
for every distinguished triangle $A \to B \to C \to A[1]$, the sequence
\[
    H(A) \longrightarrow H(B) \longrightarrow H(C)
\]
is exact. In this case, writing $H^n = H \circ [n]$, every distinguished triangle
gives a long exact sequence
\[
    \cdots \to H^n(A) \to H^n(B) \to H^n(C) \to H^{n+1}(A) \to \cdots
\]
\end{proposition}

\begin{example}\label{ex:cohomological-functors}
\begin{enumerate}
    \item The cohomology functor $H^0 \colon D(\mathcal{A}) \to \mathcal{A}$ is cohomological.
    \item For any object $X$, $\Hom_{D(\mathcal{A})}(X, -)$ is cohomological with values
    in $\Ab$.
\end{enumerate}
\end{example}


\section{Derived Functors}\label{sec:derived-functors}

\begin{definition}[Right derived functor]\label{def:right-derived-functor}
\index{derived functor!right}\index{RF@$RF$}
Let $F \colon \mathcal{A} \to \mathcal{B}$ be a left exact functor between abelian
categories, with $\mathcal{A}$ having enough injectives. The \textbf{right derived functor}
$RF \colon D^+(\mathcal{A}) \to D^+(\mathcal{B})$ is defined as follows. For
$A^\bullet \in D^+(\mathcal{A})$:
\begin{enumerate}
    \item Choose a quasi-isomorphism $A^\bullet \xrightarrow{\sim} I^\bullet$ with
    $I^\bullet \in K^+(\mathcal{I})$ (injective resolution).
    \item Define $RF(A^\bullet) = F(I^\bullet)$ (apply $F$ degreewise).
\end{enumerate}
By Theorem~\ref{thm:injective-resolution-equiv}, this is well-defined (independent
of the choice of resolution, up to canonical isomorphism in $D^+(\mathcal{B})$).

The \textbf{classical right derived functors}\index{derived functor!classical} are
\[
    R^n F(A) = H^n(RF(A)) \quad \text{for } A \in \mathcal{A}, \; n \geq 0.
\]
Note $R^0 F \cong F$ since $F$ is left exact.
\end{definition}

\begin{definition}[Left derived functor]\label{def:left-derived-functor}
\index{derived functor!left}\index{LF@$LF$}
Let $G \colon \mathcal{A} \to \mathcal{B}$ be a right exact functor with $\mathcal{A}$
having enough projectives. The \textbf{left derived functor}
$LG \colon D^-(\mathcal{A}) \to D^-(\mathcal{B})$ is defined by:
\begin{enumerate}
    \item Choose a quasi-isomorphism $P^\bullet \xrightarrow{\sim} A^\bullet$ with
    $P^\bullet \in K^-(\mathcal{P})$.
    \item Define $LG(A^\bullet) = G(P^\bullet)$.
\end{enumerate}
The classical left derived functors are $L_n G(A) = H^{-n}(LG(A))$ for $n \geq 0$,
with $L_0 G \cong G$.
\end{definition}

\begin{remark}\label{rem:derived-universal}
\index{derived functor!universal property}
The right derived functor $RF$ is characterized by the universal property: it is the
``closest approximation'' to $F$ that is an exact functor (i.e., preserves distinguished
triangles) on the derived category. More precisely, $RF$ is the right Kan extension
of $Q_{\mathcal{B}} \circ F$ along $Q_{\mathcal{A}}$.
\end{remark}

\begin{theorem}[Derived functor of a composition]\label{thm:grothendieck-spectral}
\index{Grothendieck spectral sequence}\index{spectral sequence!Grothendieck}
Let $F \colon \mathcal{A} \to \mathcal{B}$ and $G \colon \mathcal{B} \to \mathcal{C}$
be left exact functors between abelian categories with enough injectives. If $F$ sends
injective objects of $\mathcal{A}$ to $G$-acyclic objects (i.e., $R^n G(F(I)) = 0$ for
$n > 0$ and $I$ injective), then
\[
    R(G \circ F) \cong RG \circ RF.
\]
At the level of classical derived functors, there is a \textbf{Grothendieck spectral sequence}:
\[
    E_2^{p,q} = R^p G(R^q F(A)) \Longrightarrow R^{p+q}(G \circ F)(A).
\]
\end{theorem}


\section{Ext and Tor}\label{sec:ext-tor}

\begin{definition}[Ext groups]\label{def:ext}
\index{Ext}\index{Ext@$\operatorname{Ext}^n$}
Let $\mathcal{A}$ be an abelian category with enough injectives (or enough projectives).
For objects $A, B \in \mathcal{A}$, the \textbf{Ext groups} are
\[
    \operatorname{Ext}^n_{\mathcal{A}}(A, B) = R^n \Hom_{\mathcal{A}}(A, -)(B)
    = H^n(R\Hom(A, B)).
\]
Equivalently, if $I^\bullet$ is an injective resolution of $B$:
\[
    \operatorname{Ext}^n(A, B) \cong H^n(\Hom(A, I^\bullet)).
\]
If $\mathcal{A}$ has enough projectives, using a projective resolution $P_\bullet \to A$:
\[
    \operatorname{Ext}^n(A, B) \cong H^n(\Hom(P_\bullet, B)).
\]
\end{definition}

\begin{proposition}\label{prop:ext-properties}
\index{Ext!properties}
For a module category $\Mod_R$:
\begin{enumerate}
    \item $\operatorname{Ext}^0_R(M, N) \cong \Hom_R(M, N)$.
    \item $\operatorname{Ext}^n_R(M, N) = 0$ for all $n > 0$ if $M$ is projective or $N$
    is injective.
    \item The two definitions (via injective or projective resolutions) agree: this is
    the \textbf{balancing of Ext}\index{Ext!balancing}.
    \item $\operatorname{Ext}^n_R(M, N)$ classifies $n$-fold extensions of $M$ by $N$
    (Yoneda's interpretation)\index{Yoneda Ext}.
\end{enumerate}
\end{proposition}

\begin{example}\label{ex:ext-computations}
\begin{enumerate}
    \item\index{Ext!of abelian groups}
    In $\Ab$: $\operatorname{Ext}^1(\mathbb{Z}/m\mathbb{Z}, \mathbb{Z}/n\mathbb{Z})
    \cong \mathbb{Z}/\gcd(m,n)\mathbb{Z}$.
    Indeed, the projective resolution $0 \to \mathbb{Z} \xrightarrow{m} \mathbb{Z} \to \mathbb{Z}/m\mathbb{Z} \to 0$
    gives $\operatorname{Ext}^1 \cong \mathbb{Z}/n\mathbb{Z} \big/ m(\mathbb{Z}/n\mathbb{Z})
    \cong \mathbb{Z}/\gcd(m,n)\mathbb{Z}$.
    \item $\operatorname{Ext}^n_{\mathbb{Z}}(A, B) = 0$ for all $n \geq 2$ and all
    abelian groups $A, B$ (since $\mathbb{Z}$ has global dimension 1).
\end{enumerate}
\end{example}

\begin{definition}[Tor groups]\label{def:tor}
\index{Tor}\index{Tor@$\operatorname{Tor}_n$}
For a ring $R$ and modules $M \in \Mod_{R^{\op}}$, $N \in \Mod_R$, the \textbf{Tor groups} are
\[
    \operatorname{Tor}_n^R(M, N) = L_n(M \otimes_R -)(N) = H^{-n}(M \otimes_R^L N).
\]
If $P_\bullet \to N$ is a projective resolution:
\[
    \operatorname{Tor}_n^R(M, N) \cong H_n(M \otimes_R P_\bullet).
\]
\end{definition}

\begin{proposition}\label{prop:tor-properties}
\index{Tor!properties}
\begin{enumerate}
    \item $\operatorname{Tor}_0^R(M, N) \cong M \otimes_R N$.
    \item $\operatorname{Tor}_n^R(M, N) = 0$ for all $n > 0$ if $M$ or $N$ is flat.
    \item Tor is balanced: $L_n(M \otimes_R -)(N) \cong L_n(- \otimes_R N)(M)$.
    \item $M$ is flat if and only if $\operatorname{Tor}_1^R(M, N) = 0$ for all $N$.
\end{enumerate}
\end{proposition}

\begin{example}\label{ex:tor-computations}
\begin{enumerate}
    \item\index{Tor!of abelian groups}
    $\operatorname{Tor}_1^{\mathbb{Z}}(\mathbb{Z}/m\mathbb{Z}, \mathbb{Z}/n\mathbb{Z})
    \cong \mathbb{Z}/\gcd(m,n)\mathbb{Z}$.
    Using $0 \to \mathbb{Z} \xrightarrow{m} \mathbb{Z} \to \mathbb{Z}/m\mathbb{Z} \to 0$:
    $\operatorname{Tor}_1 = \Ker(\mathbb{Z}/n\mathbb{Z} \xrightarrow{m} \mathbb{Z}/n\mathbb{Z})
    \cong \mathbb{Z}/\gcd(m,n)\mathbb{Z}$.
    \item $\operatorname{Tor}_n^{\mathbb{Z}}(A, B) = 0$ for $n \geq 2$.
\end{enumerate}
\end{example}


\section{Long Exact Sequences}\label{sec:long-exact-derived}

\begin{theorem}\label{thm:les-derived}
\index{long exact sequence!from derived functors}
Let $F \colon \mathcal{A} \to \mathcal{B}$ be a left exact functor and
$0 \to A' \to A \to A'' \to 0$ a short exact sequence in $\mathcal{A}$.
Then there is a long exact sequence:
\[
    0 \to FA' \to FA \to FA'' \xrightarrow{\delta^0} R^1 FA' \to R^1 FA \to R^1 FA''
    \xrightarrow{\delta^1} R^2 FA' \to \cdots
\]
Dually, for a right exact functor $G$ and $0 \to A' \to A \to A'' \to 0$:
\[
    \cdots \to L_2 GA'' \xrightarrow{\delta_2} L_1 GA' \to L_1 GA \to L_1 GA''
    \xrightarrow{\delta_1} GA' \to GA \to GA'' \to 0.
\]
\end{theorem}

\begin{proof}
By Proposition~\ref{prop:ses-to-triangle}, the short exact sequence gives a distinguished
triangle $A' \to A \to A'' \to A'[1]$ in $D^+(\mathcal{A})$. Applying $RF$ (which
preserves distinguished triangles) gives a distinguished triangle
\[
    RF(A') \to RF(A) \to RF(A'') \to RF(A')[1]
\]
in $D^+(\mathcal{B})$. The cohomological functor $H^0$ (Proposition~\ref{prop:cohomological-functor})
then yields the long exact sequence, since $H^n(RF(A)) = R^n F(A)$.
\end{proof}

\begin{corollary}\label{cor:ext-les}
\index{Ext!long exact sequence}
For any short exact sequence $0 \to A' \to A \to A'' \to 0$ in $\Mod_R$ and any module $M$:
\begin{enumerate}
    \item (Covariant) There is a long exact sequence
    \[
        0 \to \Hom(M, A') \to \Hom(M, A) \to \Hom(M, A'') \to
        \operatorname{Ext}^1(M, A') \to \cdots
    \]
    \item (Contravariant) There is a long exact sequence
    \[
        0 \to \Hom(A'', M) \to \Hom(A, M) \to \Hom(A', M) \to
        \operatorname{Ext}^1(A'', M) \to \cdots
    \]
\end{enumerate}
\end{corollary}

\begin{corollary}\label{cor:tor-les}
\index{Tor!long exact sequence}
For any short exact sequence $0 \to A' \to A \to A'' \to 0$ in $\Mod_R$ and any module $M$:
\[
    \cdots \to \operatorname{Tor}_1(M, A') \to \operatorname{Tor}_1(M, A) \to
    \operatorname{Tor}_1(M, A'') \to M \otimes A' \to M \otimes A \to M \otimes A'' \to 0.
\]
\end{corollary}

\begin{example}\label{ex:les-application}
\index{long exact sequence!computation example}
Consider the exact sequence $0 \to \mathbb{Z} \xrightarrow{n} \mathbb{Z} \to \mathbb{Z}/n\mathbb{Z} \to 0$
in $\Ab$. Applying $\Hom(-, A)$ for an abelian group $A$:
\[
    0 \to \Hom(\mathbb{Z}/n, A) \to \Hom(\mathbb{Z}, A) \xrightarrow{n}
    \Hom(\mathbb{Z}, A) \to \operatorname{Ext}^1(\mathbb{Z}/n, A) \to 0.
\]
This gives $\Hom(\mathbb{Z}/n, A) \cong A[n] = \{a \in A : na = 0\}$ and
$\operatorname{Ext}^1(\mathbb{Z}/n, A) \cong A/nA$.
\end{example}


\section{Derived Category in Practice}\label{sec:derived-practice}

\begin{definition}[RHom and derived tensor]\label{def:rhom-ltensor}
\index{RHom}\index{derived tensor product}\index{tensor product!derived}
For $\mathcal{A} = \Mod_R$:
\begin{enumerate}
    \item $R\Hom_R(M, N) \in D^+(\Ab)$ is the right derived functor of $\Hom_R(M, -)$.
    Its cohomology gives $H^n(R\Hom(M, N)) = \operatorname{Ext}^n_R(M, N)$.
    \item $M \otimes_R^L N \in D^-(\Ab)$ is the left derived functor of $M \otimes_R -$.
    Its cohomology gives $H^{-n}(M \otimes_R^L N) = \operatorname{Tor}_n^R(M, N)$.
\end{enumerate}
These extend to functors on derived categories:
\[
    R\Hom \colon D^-(\Mod_R)^{\op} \times D^+(\Mod_R) \to D^+(\Ab),
\]
\[
    \otimes_R^L \colon D^-(\Mod_{R^{\op}}) \times D^-(\Mod_R) \to D^-(\Ab).
\]
\end{definition}

\begin{proposition}[Derived Hom--Tensor adjunction]\label{prop:derived-adjunction}
\index{adjunction!derived Hom--Tensor}
There is a natural isomorphism in $D(\Ab)$:
\[
    R\Hom_R(M \otimes_S^L N, \; P) \cong R\Hom_S(N, \; R\Hom_R(M, P))
\]
for appropriate bimodule structures.
\end{proposition}

\begin{proposition}[Derived category of a hereditary category]\label{prop:hereditary-derived}
\index{hereditary category}\index{derived category!hereditary case}
Let $\mathcal{A}$ be a hereditary abelian category (i.e., $\operatorname{Ext}^n = 0$ for $n \geq 2$).
Then every object of $D^b(\mathcal{A})$ is isomorphic to the direct sum of its cohomology
objects (each placed in its own degree):
\[
    A^\bullet \cong \bigoplus_{n \in \mathbb{Z}} H^n(A^\bullet)[-n] \quad \text{in } D^b(\mathcal{A}).
\]
\end{proposition}

\begin{remark}\label{rem:derived-cat-philosophy}
\index{derived category!philosophy}
The philosophy of derived categories can be summarized as follows:
\begin{enumerate}
    \item Objects of $\mathcal{A}$ are replaced by complexes, which carry ``higher-order''
    information through their cohomology in all degrees.
    \item Functors that are only partially exact (left or right exact) become exact on derived
    categories, at the cost of working with complexes.
    \item The derived category provides a unified framework for cohomological computations:
    spectral sequences, universal coefficient theorems, K\"unneth formulas, etc.\
    all arise from the triangulated structure.
\end{enumerate}
\end{remark}

\begin{notation}\label{not:derived-summary}
\index{derived category!notation summary}
We summarize the key notation:
\begin{center}
\renewcommand{\arraystretch}{1.3}
\begin{tabular}{ll}
\hline
$\operatorname{Ch}(\mathcal{A})$ & Category of chain complexes \\
$K(\mathcal{A})$ & Homotopy category \\
$D(\mathcal{A})$ & Derived category \\
$D^+, D^-, D^b$ & Bounded variants \\
$RF$ & Right derived functor \\
$LG$ & Left derived functor \\
$R\Hom$ & Derived Hom \\
$\otimes^L$ & Derived tensor product \\
$\operatorname{Ext}^n(A,B)$ & $= H^n(R\Hom(A,B))$ \\
$\operatorname{Tor}_n(M,N)$ & $= H^{-n}(M \otimes^L N)$ \\
\hline
\end{tabular}
\end{center}
\end{notation}


\section{Exercises}\label{sec:derived-exercises}

\begin{exercise}\label{exo:acyclic-cone}
\index{mapping cone!and quasi-isomorphism}
Let $f \colon A^\bullet \to B^\bullet$ be a chain map. Show that $f$ is a quasi-isomorphism
if and only if $\operatorname{Cone}(f)$ is acyclic.
\end{exercise}

\begin{exercise}\label{exo:homotopy-category-additive}
Show that the homotopy category $K(\mathcal{A})$ is additive. Verify that it is not
abelian by finding an explicit example (e.g., in $K(\Ab)$) of a morphism that has no
kernel in $K(\Ab)$.
\end{exercise}

\begin{exercise}\label{exo:shift-exact}
\index{shift functor!exactness}
Show that the shift functor $[1] \colon D(\mathcal{A}) \to D(\mathcal{A})$ is an
exact functor of triangulated categories (sends distinguished triangles to
distinguished triangles).
\end{exercise}

\begin{exercise}\label{exo:ext-computation}
\index{Ext!computation}
Let $k$ be a field. Compute $\operatorname{Ext}^n_{k[x]}(k, k)$ for all $n \geq 0$,
where $k$ is viewed as a $k[x]$-module via $x \mapsto 0$. (Use the projective resolution
$0 \to k[x] \xrightarrow{x} k[x] \to k \to 0$.)
\end{exercise}

\begin{exercise}\label{exo:tor-flat}
\index{Tor!and flatness}
Let $R$ be a commutative ring and $M$ an $R$-module.
\begin{enumerate}
    \item Show that $M$ is flat if and only if $\operatorname{Tor}_1^R(M, R/I) = 0$ for
    every ideal $I \subseteq R$.
    \item Compute $\operatorname{Tor}_n^{\mathbb{Z}}(\mathbb{Q}, \mathbb{Z}/m\mathbb{Z})$
    for all $n, m$.
\end{enumerate}
\end{exercise}

\begin{exercise}\label{exo:derived-ses}
\index{derived functor!long exact sequence}
Let $0 \to M' \to M \to M'' \to 0$ be an exact sequence of $R$-modules and $N$ any
$R$-module. Write down explicitly the long exact sequences in $\operatorname{Ext}$ and
$\operatorname{Tor}$ obtained by applying $\Hom_R(N, -)$, $\Hom_R(-, N)$, and
$N \otimes_R -$.
\end{exercise}

\begin{exercise}\label{exo:global-dimension}
\index{global dimension}\index{projective dimension}
The \textbf{projective dimension} $\operatorname{pd}(M)$ of a module $M$ is the infimum of
lengths of projective resolutions of $M$. The \textbf{global dimension} of $R$ is
$\operatorname{gldim}(R) = \sup_M \operatorname{pd}(M)$.
\begin{enumerate}
    \item Show that $\operatorname{pd}(M) \leq n$ if and only if $\operatorname{Ext}^{n+1}_R(M, N) = 0$
    for all $N$.
    \item Show that $\operatorname{gldim}(\mathbb{Z}) = 1$.
    \item Show that $\operatorname{gldim}(k[x_1, \ldots, x_n]) = n$ for a field $k$
    (this is the Hilbert syzygy theorem\index{Hilbert syzygy theorem}).
\end{enumerate}
\end{exercise}

\begin{exercise}\label{exo:triangulated-hom}
\index{triangulated category!Hom long exact sequence}
Let $\mathcal{T}$ be a triangulated category and $A \xrightarrow{f} B \xrightarrow{g} C
\xrightarrow{h} A[1]$ a distinguished triangle. For any object $X$, show there are long
exact sequences:
\begin{align*}
    \cdots &\to \Hom(X, A) \to \Hom(X, B) \to \Hom(X, C) \to \Hom(X, A[1]) \to \cdots \\
    \cdots &\to \Hom(C, X) \to \Hom(B, X) \to \Hom(A, X) \to \Hom(C[-1], X) \to \cdots
\end{align*}
\end{exercise}

\begin{exercise}\label{exo:derived-cat-Hom-Ext}
\index{derived category!Hom and Ext}
Show that in $D(\mathcal{A})$, for objects $A, B \in \mathcal{A}$ (viewed as complexes
concentrated in degree 0):
\[
    \Hom_{D(\mathcal{A})}(A, B[n]) \cong \operatorname{Ext}^n_{\mathcal{A}}(A, B)
    \quad \text{for all } n \geq 0.
\]
This is one of the key motivations for derived categories: the Ext groups are simply
Hom spaces with shifts.
\end{exercise}

\begin{exercise}\label{exo:verdier-quotient}
\index{Verdier quotient}\index{triangulated category!Verdier quotient}
(Verdier quotient.) Let $\mathcal{T}$ be a triangulated category and
$\mathcal{N} \subseteq \mathcal{T}$ a \textbf{thick subcategory} (a full triangulated
subcategory closed under direct summands). Define the Verdier quotient $\mathcal{T}/\mathcal{N}$
as the localization at morphisms whose cone lies in $\mathcal{N}$.
\begin{enumerate}
    \item Show that $D(\mathcal{A}) \cong K(\mathcal{A}) / \operatorname{Ac}(\mathcal{A})$,
    where $\operatorname{Ac}(\mathcal{A})$ is the thick subcategory of acyclic complexes.
    \item Show that the Verdier quotient inherits a triangulated structure.
\end{enumerate}
\end{exercise}
